
The dominant explanation for why neural networks rely on spurious features---that such features produce stronger gradient signals during training---is empirically incorrect. We set out to characterize when deep networks learn spurious correlations, expecting to find that simple, spurious patterns commandeer gradient flow early in training. Instead, we discovered that minority-group samples---those requiring invariant feature learning---produce 5.7 times higher gradient norms than spurious-aligned samples. The gradient competition hypothesis, implicitly assumed across much of the robustness literature, is inverted in our measurements.

This finding has implications for robustness interventions. Methods that attempt to rebalance gradient contributions between easy and hard samples assume that easy (spurious-aligned) samples dominate the learning signal. Our measurements reveal the opposite: spurious-aligned samples achieve low loss quickly and subsequently contribute less to parameter updates, while minority samples continue generating high gradients precisely because they remain unsolved.

\subsubsection{1.1 The Spurious Correlation Problem}

Deep neural networks trained with empirical risk minimization (ERM) exploit statistical shortcuts present in training data. On the Waterbirds benchmark, models learn to classify birds by background habitat rather than bird morphology---a spurious correlation that holds for most training examples but fails on minority groups where the correlation breaks. Despite achieving high average accuracy, such models may correctly classify fewer than 70\% of minority-group samples.

The field has developed various interventions: Group DRO minimizes worst-group loss when group annotations are available; JTT identifies likely spurious predictions via early training errors and upweights them in a second training phase; SAM seeks flat minima that may generalize better across distribution shifts.

\subsubsection{1.2 Beyond Detection: Understanding the Mechanism}

Most work focuses on detecting or mitigating spurious reliance post-hoc, rather than understanding why networks preferentially learn spurious features. The simplicity bias hypothesis provides a partial answer---simpler features are learned first---but the mechanism by which this bias operates during gradient-based optimization remains underspecified.

A natural assumption is that spurious features, being simpler, produce stronger gradient signals that accelerate their learning. This ''gradient competition'' view suggests that spurious and invariant features compete for representational capacity, with spurious features winning through gradient magnitude. If true, interventions could target this competition directly.

\subsubsection{1.3 Our Investigation}

We test this mechanistic hypothesis through two complementary experiments:

\textbf{Existence (H-E1):} We first verify that spurious feature dominance occurs early in training, measuring the ratio of GradCAM attribution on spurious (background) versus core (bird) regions across 50 training epochs.

\textbf{Mechanism (H-M1):} We then measure gradient norm magnitudes for spurious-aligned versus minority samples, testing whether spurious features receive stronger optimization signals.

Our results confirm spurious dominance from epoch 0 (attribution ratio 1.35) but falsify the gradient competition mechanism. The expected gradient ratio of >1.5 favoring spurious features inverts to 0.18---minority samples generate far stronger gradients than spurious-aligned ones.

\subsubsection{1.4 Contributions}

This paper makes three contributions:

\begin{enumerate}
\item \textbf{Empirical falsification} of the gradient competition hypothesis for spurious feature learning, demonstrating that gradient norm ratios are inverted from theoretical predictions (observed 0.18 vs. expected >1.5).
\end{enumerate}

\begin{enumerate}
\item \textbf{Confirmation and characterization} of simplicity bias dynamics, showing spurious attribution dominates from initialization with ratio 1.35, peaking at 1.59 around epoch 13.
\end{enumerate}

\begin{enumerate}
\item \textbf{Mechanistic reframing} proposing that spurious dominance arises from convergence speed rather than gradient magnitude---spurious patterns occupy simpler loss landscape regions that achieve low loss faster.
\end{enumerate}
